\par\bigskip\pdfbookmark[0]{Zur Theorie der biquadratischen Reste. IV.}{nachlass-biquadratic-IV}\centerline{\textbf{[IV.]}}\medskip
\renewcommand{\Article}[1]{\par\medskip\Needspace{4\baselineskip}\pdfbookmark[1]{Art. #1}{nachlass-biquadratic-IV-art-#1}\centerline{\textbf{#1.}}\nobreak\smallskip}
\Article{1}
\[
\begin{array}{lll}
\text{Biquadratischer Rest?}&m=a+bi;&aa+bb=d\\[3pt]
\text{Modulus}&M=A+Bi,&AA+BB=D
\end{array}
\]
\[
\frac{mD}{M}=\mu\qquad
\mu=\alpha+\beta i,\quad\alpha=aA+bB,\quad\beta=Ab-Ba
\]
\[
\xi+\eta i=\pi;\quad\pi m=x+yi=p;\quad\pi M=X+Yi=P
\]
Relationen
\[
\begin{array}{lll}
x=a\xi-b\eta&d\xi=ax+by&D\xi=AX+BY\\
y=b\xi+a\eta&d\eta=-bx+ay&D\eta=-BX+AY\\
X=A\xi-B\eta&dX=\alpha x+\beta y&Dx=\alpha X-\beta Y\\
Y=B\xi+A\eta&dY=-\beta x+\alpha y&Dy=\beta X+\alpha Y
\end{array}
\]

\SourceLeaf{353}
\[
\begin{aligned}
\beta\xi&=-Bx+bX=Ay-aY\\
\beta\eta&=-Ax+aX=-By+bY\\
\alpha\xi&=Ax+aY=By+aX\\
\alpha\eta&=-Bx+aY=Ay-bX
\end{aligned}
\]

Diejenigen $\pi$, wo $\xi$ und $\eta$ zwischen $0$ und $\frac12$ liegen, sollen durch $\pi^0$ bezeichnet werden, und die entsprechenden $p$ und $P$ durch $p^0$ und $P^0$; diejenigen $\pi$, wo $\eta=0$ und $\xi$ zwischen $0$ und $\frac12$, durch $\pi'$; die, wo $\xi=\frac12$ und $\eta$ zwischen $0$ und $\frac12$, durch $\pi''$; diejenigen $\pi$, wo $\eta=\frac12$ und $\xi$ zwischen $0$ und $\frac12$, durch $\pi'''$; endlich die wo $\xi=0$ und $\eta$ zwischen $0$ und $\frac12$, durch $\pi''''$.

Der Decident von $\frac{m}{M}$ wird so gefunden:

Man sammle alle \emph{ganzen} $P^0$, für welche mithin $x^0$ und $y^0$ gebrochen sein werden; die respectiven Intensoren von $p^0$ seien $t^0$ d.\,i. die Zahlen $0,1.\ 2,3$, je nachdem
\[
\begin{array}{lllll}
[x^0]&\text{gerade,}&\text{ungerade,}&\text{ungerade,}&\text{gerade}\\
\relax[y^0]&\text{gerade,}&\text{gerade,}&\text{ungerade,}&\text{ungerade}
\end{array}
\]
So ist der gesuchte Decident $=\Sigma\pm t^0$, wo das obere Zeichen für gerade $P^0$, das untere für die ungeraden zu nehmen ist.

Dies ist die \emph{erste} Methode.

\Article{2}
Wir wollen nun die einzelnen $P^0$ nach den Werthen von $Y^0$ zusammenordnen. Indem wir uns auf den Fall einschränken, wo $a,b,A,B$ positiv sind, ist der kleinste Werth von $Y^0\dotstwo+1$, der grösste $\frac12(A+B-1)$. Für jeden bestimmten Werth von $Y^0$ müssen die Werthe von $X^0$ zwischen bestimmten Grenzen liegen, nemlich

I. wenn $A-B$ positiv ist
\[
\begin{array}{c|c|c}
\text{wenn}&\text{zwischen}&\text{und}\\
Y<\frac12B&-\frac{BY^0}{A}&\frac{AY^0}{B}\\[5pt]
Y=\frac12B&-\frac{BB}{2A}&\frac12A\\[5pt]
Y>\frac12B\text{ und }<\frac12A&-\frac{BY^0}{A}&-\frac{BY^0}{A}+\frac{D}{2A}\\[5pt]
Y>\frac12A&\frac{AY^0}{B}-\frac{D}{2B^0}&-\frac{BY^0}{A}+\frac{D}{2A}
\end{array}
\]

\SourceLeaf{354}
II. Wenn $A-B$ negativ ist,
\[
\begin{array}{c|c|c}
\text{wenn}&\text{zwischen}&\text{und}\\
Y<\frac12A&-\frac{BY^0}{A}&\frac{AY^0}{B}\\[5pt]
Y>\frac12A\text{ und }<\frac12B&\frac{AY^0}{B}-\frac{D}{2B}&\frac{AY^0}{B}\\[5pt]
Y=\frac12B&\frac{AB-D}{2B}&\frac12A\\[5pt]
Y>\frac12B&\frac{AY^0}{B}-\frac{D}{2B}&-\frac{BY^0}{A}+\frac{D}{2A}
\end{array}
\]

In den kleinern Grenzen ist entweder $\xi=0$ oder $\eta=\frac12$, in den grössern Grenzen hingegen entweder $\eta=0$ oder $\xi=\frac12$. Es lässt sich leicht beweisen, dass nie die Grenzen von $x$ ganze Zahlen sind.

\Article{3}
Wir wollen nun die Partialsummen für jedes bestimmte $Y^0$ auf eine andere Weise darstellen. Auf den Grenzen wird $p$ bestimmte Werthe haben, die durch $p^*,p^{**}$ bezeichnet werden mögen, und während $X$ stetig von der einen Grenze zur andern sich ändert, wird $p$ stetig von $p^*$ zu $p^{**}$ übergehen. Allein die in $[p]$ enthaltene ganze Zahl wird hiebei sprungsweise geändert, indem immer entweder der reelle oder der imaginäre Theil sich um Eine Einheit ändert. Es geschehen die Aenderungen bei den Werthen von $X$
\[
X',\ X'',\ X'''\ldots X^\mu
\]
die bereits nach ihrer Grösse geordnet sind und denen die Werthe von $p$
\[
p',\ p''.\ p'''\ldots
\]
entsprechen.

Das letzte $X^\mu$ kann auch mit $X^{**}$ identisch sein, wenn $\beta$ positiv, oder $X'$ mit $X^*$ identisch etc.

Die $x$ sind hier zunehmend, also wenn $x^{**}$ eine ganze Zahl, wird sie für $x^\mu$ gezählt.

Die $y$ sind zunehmend bei positiven $\beta$, da wird $y^{**}$ ganz mitgezählt
\begin{center}abnehmend bei negativen $\beta$, da wird $y^*$ mitgezählt.\end{center}

\SourceLeaf{355}
Die Intensoren von $p^*,p^{**}$ seien $\lambda^*$ und $\lambda^{**}$
\[
\begin{array}{rll}
\text{der Intensor von }p'&\text{an bis }p''\dotstwo&\lambda^*+\delta'\\
p''&\text{bis }p'''\dotstwo&\lambda^*+\delta'+\delta''\\
p^\mu&\text{bis }p^{**}\dotstwo&\lambda^*+\delta'+\delta''\dotstwo+\delta^\mu=\lambda^{**}
\end{array}
\]
so wird die Partialsumme, in sofern $Y^0$ gerade,
\[
\begin{aligned}
={}&\lambda^*(g-h)+(\lambda^*+\delta')(g'-h')\\
&+(\lambda^*+\delta'+\delta'')(g''-h'')+\text{etc.}+\lambda^{**}(g^\mu-h^\mu)
\end{aligned}
\]
wo $g$ die Anzahl der geraden $X^0$ von $X^*$ bis $X'$, $h$ die der ungeraden bedeutet.

\Article{4}
Diese Formel lässt sich auch so darstellen:
\[
\begin{aligned}
&\lambda^*\{[\tfrac12X^*-\tfrac12]-[\tfrac12X^*]\}\\
&+\delta'\{[\tfrac12X'-\tfrac12]-[\tfrac12X']\}\\
&+\delta''\{[\tfrac12X''-\tfrac12]-[\tfrac12X'']\}\\
&+\text{etc.}\\
&+\delta^\mu\{[\tfrac12X^\mu-\tfrac12]-[\tfrac12X^\mu]\}\\
&-\lambda^{**}\{[\tfrac12X^{**}-\tfrac12]-[\tfrac12X^{**}]\}
\end{aligned}
\]
oder durch
\[
\lambda^*\varepsilon^*+\delta'\varepsilon'+\delta''\varepsilon''+\text{etc.}+\delta^\mu\varepsilon^\mu-\lambda^{**}\varepsilon^{**}
\]
\[
\begin{array}{rll}
\text{wo allgemein }\varepsilon=0&\text{wenn }[X]\text{ ungerade}\\
\text{und }=-1&\text{wenn }[X]\text{ gerade ist und }Y\text{ gerade}\\
+1&\text{wenn }[X]\text{ gerade und }Y\text{ ungerade.}
\end{array}
\]

Für $\delta$ hingegen hat man die Werthe
\[
\begin{array}{ll|rr}
&&\multicolumn{2}{c}{\beta}\\
&&\text{positiv}&\text{negativ}\\
\text{wenn }x\text{ eine ganze gerade Zahl,}&[y]\text{ gerade}&-1&-1\\
&[y]\text{ ungerade}&+1&+1\\
x\text{ eine ganze ungerade Zahl,}&[y]\text{ gerade}&+1&+1\\
&[y]\text{ ungerade}&-1&-1\\
y\text{ eine ganze gerade Zahl,}&[x]\text{ gerade}&-3&+3\\
&[x]\text{ ungerade}&-1&+1\\
y\text{ eine ganze ungerade Zahl,}&[x]\text{ gerade}&+3&-3\\
&[x]\text{ ungerade}&+1&-1
\end{array}
\]

\SourceLeaf{356}
\Article{5}
Hieraus leiten wir folgende zweite Methode den Decidenten zu bestimmen ab.

I. Man sammle alle Combinationen von ganzen Werthen von $Y$ und $x$, die folgende Eigenschaften haben
\begin{enumerate}
\item dass $\xi=\frac{Ax+bY}{\alpha}$ zwischen $0$ und $\frac12$ falle, wobei die zweite Grenze inclusive genommen wird
\item dass $\eta=\frac{-Bx+aY}{\alpha}$ zwischen $0$ und $\frac12$ falle, die erste Grenze inclusive genommen.
\end{enumerate}

II. Man berechne dafür
\[
\begin{aligned}
X&=\frac{Dx+\beta Y}{\alpha}\\[3pt]
y&=\frac{\beta x+dY}{\alpha}
\end{aligned}
\]

III. Man lasse alle diejenigen weg, wo $[X]$ eine ungerade Zahl ist, und theile die übrigen, wo $[X]$ gerade ist, in zwei Classen;

in die erste Classe setze man diejenigen, wo zugleich
\[
\begin{array}{l}
Y\text{ gerade, }x\text{ gerade, }[y]\text{ gerade}\\
\text{oder wo eine dieser Bedingungen Statt findet;}
\end{array}
\]
in die zweite Classe diejenigen, wo zwei dieser Bedingungen oder gar keine Statt hat,
\[
\begin{array}{rl}
\text{oder in I.}&\text{wo }[Y+x+y]\text{ gerade}\\
\text{II.}&\text{wo }[Y+x+y]\text{ ungerade}
\end{array}
\]
und nenne den Ueberschuss der Anzahl in der ersten Classe über die in der zweiten $c$.

IV. Man sammle alle Combinationen von ganzen Werthen von $Y$ und $y$, die folgende Eigenschaften haben:
\begin{enumerate}
\item dass $\xi=\frac{Ay-aY}{\beta}$ zwischen $0$ und $\frac12$ falle, die erste Grenze inclusive, wenn $\beta$ negativ, die zweite inclusive, wenn $\beta$ positiv;
\item dass $\eta=\frac{-By+bY}{\beta}$ zwischen $0$ und $\frac12$ falle, die erste Grenze inclusive bei positivem $\beta$, die zweite bei negativem.
\end{enumerate}
\BookRule{.3\textwidth}
